% ============================================================
%  Seizure Prediction – IEEE LaTeX Report (Overleaf-ready)
%  IEEEtran two-column conference style
% ============================================================
\documentclass[conference]{IEEEtran}

% ── Packages ─────────────────────────────────────────────
\usepackage[T1]{fontenc}
\usepackage[utf8]{inputenc}
\usepackage{amsmath,amssymb,amsfonts}
\usepackage{graphicx}
\usepackage{booktabs}
\usepackage{tabularx}
\usepackage{multirow}
\usepackage{xcolor}
\usepackage{hyperref}
\usepackage{cite}
\usepackage{balance}
\usepackage{microtype}
\usepackage{float}
\usepackage{caption}
\usepackage{subcaption}
\usepackage{algorithm}
\usepackage{algpseudocode}
\usepackage{siunitx}

% ── Colour palette ───────────────────────────────────────
\definecolor{navyblue}{HTML}{1A2E4A}
\definecolor{tealgreen}{HTML}{0A9396}
\definecolor{coralred}{HTML}{EE6C4D}
\definecolor{goldyellow}{HTML}{E9C46A}
\definecolor{lightgray}{HTML}{F4F4F4}

% ── Hyperref setup ───────────────────────────────────────
\hypersetup{
  colorlinks = true,
  linkcolor  = navyblue,
  citecolor  = tealgreen,
  urlcolor   = tealgreen
}

% ── Caption style ────────────────────────────────────────
\captionsetup{font=small, labelfont=bf}

% ── Graphics path ────────────────────────────────────────
\graphicspath{{figures/}}

% ============================================================
\begin{document}

% ── Title ────────────────────────────────────────────────
\title{Seizure Prediction with Logistic Regression:\\
  The Effect of Preprocessing Order, Model Complexity,\\
  and Regularisation on Generalisation}

\author{
  \IEEEauthorblockN{[Student Name]}
  \IEEEauthorblockA{
    \textit{Department of Computer Science / Biomedical Engineering}\\
    \textit{[University Name], [City, Country]}\\
    student.email@university.edu
  }
}

\maketitle

% ── Abstract ─────────────────────────────────────────────
\begin{abstract}\sloppy
This paper investigates how preprocessing pipeline design, regularisation
strategy, and class-imbalance handling interact to determine generalisation
performance in epileptic seizure prediction. Three simulated EEG-derived
feature datasets---modelled after CHB-MIT ($n{=}5{,}000$, $9.2\%$ seizure),
Bonn-EEG ($n{=}2{,}000$, $21.1\%$), and a Kaggle-style corpus
($n{=}1{,}000$, $38.7\%$)---serve as benchmarks with realistic separability
and $3\%$ label noise. Two preprocessing pipelines are compared: Pipeline~A
(\texttt{StandardScaler}~$\!\to\!$~\texttt{SelectKBest} ANOVA-F) and
Pipeline~B (\texttt{MinMaxScaler}$\to$\texttt{StandardScaler}$\to$PCA).
L1, L2, and Elastic Net regularised logistic regression models are trained
across an 11-point $C$ grid. Overfitting and underfitting are deliberately
demonstrated via validation and learning curves. Three class-imbalance
correction techniques (SMOTE, random undersampling, class weighting) are
evaluated using PR-AUC---the preferred metric for imbalanced clinical data.
Results show that: (1)~preprocessing order significantly affects PR-AUC
($+0.213$ on CHB-MIT); (2)~L2 achieves the best mean PR-AUC across the $C$
sweep; (3)~Elastic Net outperforms pure L1 and converges to L2 after
upstream feature selection; and (4)~the imbalance $\times$ regulariser
interaction is strongest on severely imbalanced data.
\end{abstract}

\begin{IEEEkeywords}
epileptic seizure prediction, logistic regression, preprocessing pipelines,
regularisation, class imbalance, SMOTE, PR-AUC
\end{IEEEkeywords}

% ============================================================
\section{Introduction}
\label{sec:intro}

Epilepsy affects approximately 50~million people worldwide, making seizure
prediction one of the highest-impact problems in clinical machine
learning~\cite{WHO2019}. The goal is to identify ictal (seizure) epochs in
EEG recordings before or at onset, enabling timely intervention. The task is
characterised by three interacting challenges: \emph{severe class imbalance}
(seizure epochs represent only 2--10\% of recording time), \emph{high-dimensional
correlated feature spaces} arising from multi-channel EEG, and
\emph{sensitivity to preprocessing artefacts} such as baseline drift and
muscle noise.

Despite the abundance of deep-learning approaches, logistic regression
remains a powerful and interpretable baseline. Its probabilistic output
facilitates calibrated risk thresholding, and its regularisation parameter
$C$ (inverse regularisation strength) provides a clean knob for studying the
bias-variance trade-off. This report systematically investigates how the
\emph{ordering} of preprocessing steps, the \emph{choice of regulariser}
(L1, L2, Elastic Net), and the \emph{method of class-imbalance correction}
interact to determine out-of-sample generalisation, measured primarily by
precision-recall area under the curve (PR-AUC)---the metric of choice for
imbalanced biomedical datasets~\cite{Davis2006}.

The remainder of this paper is structured as follows.
Section~\ref{sec:datasets} describes the three benchmark datasets.
Section~\ref{sec:preprocessing} details the two preprocessing pipelines.
Sections~\ref{sec:baseline}--\ref{sec:imbalance} present baseline results,
the overfitting/underfitting study, the regularisation comparison, and
imbalance handling experiments.
Section~\ref{sec:analysis} synthesises the findings.
Section~\ref{sec:conclusion} concludes.

% ============================================================
\section{Dataset Collection}
\label{sec:datasets}

Three datasets are used in this study. Each mirrors the statistical
properties of a well-known public epilepsy corpus while remaining fully
reproducible without data-access restrictions. Datasets are generated via
scikit-learn's \texttt{make\_classification}~\cite{Pedregosa2011} with
controlled class separability (\texttt{class\_sep}), $3\%$ label noise
(\texttt{flip\_y=0.03}), and correlated redundant features mimicking
EEG spectral band overlap (delta, theta, alpha, beta).

\begin{table}[h]
\centering
\caption{Dataset Summary}
\label{tab:datasets}
\begin{tabular}{lrrrrr}
\toprule
\textbf{Dataset} & $\mathbf{n}$ & \textbf{Seiz.\%} &
\textbf{Feats} & \textbf{Sep.} & \textbf{Analogue} \\
\midrule
CHB-MIT  & 5{,}000 & 9.2\%  & 64 & 0.9 & CHB-MIT EEG~\cite{Goldberger2010} \\
Bonn-EEG & 2{,}000 & 21.1\% & 32 & 1.1 & Bonn Univ.~\cite{Andrzejak2001}  \\
Kaggle   & 1{,}000 & 38.7\% & 20 & 1.3 & Kaggle~\cite{Kaggle}             \\
\bottomrule
\end{tabular}
\end{table}

\textbf{CHB-MIT} mirrors the paediatric scalp EEG dataset~\cite{Goldberger2010},
characterised by long inter-ictal segments and severe imbalance.
\textbf{Bonn-EEG} approximates the widely-cited intracranial
depth-electrode recordings~\cite{Andrzejak2001}, with moderate imbalance and
32 spectral features.
\textbf{Kaggle}~\cite{Kaggle} contains near-balanced, low-dimensional
extracted features, serving as an upper-bound reference for attainable
performance.

\begin{figure}[htbp]
  \centering
  \includegraphics[width=\columnwidth]{fig1_dataset_distribution.png}
  \caption{Class distribution across the three benchmark datasets.
           CHB-MIT exhibits the most severe imbalance (9.2\% seizure epochs).}
  \label{fig:datasets}
\end{figure}

% ============================================================
\section{Preprocessing Pipelines}
\label{sec:preprocessing}

Students must design at least two different preprocessing pipelines whose
ordering encodes a specific research insight~\cite{Hastie2009}.

\subsection{Pipeline A: Normalisation $\to$ Feature Selection}

\begin{algorithm}[h]
\caption{Pipeline A}
\label{alg:pipeA}
\begin{algorithmic}[1]
  \Require Raw feature matrix $\mathbf{X} \in \mathbb{R}^{m \times d}$
  \State $\mathbf{X}' \gets \texttt{StandardScaler}(\mathbf{X})$
    \Comment{$\mu{=}0$, $\sigma{=}1$ per feature}
  \State $\mathbf{X}'' \gets \texttt{SelectKBest}(\mathbf{X}', f_\text{ANOVA}, k{=}15)$
    \Comment{univariate ANOVA-F ranking}
  \State \Return $\mathbf{X}''$
\end{algorithmic}
\end{algorithm}

\noindent\textbf{Key insight:} Normalisation must precede scoring. Raw EEG
power features span several orders of magnitude across frequency bands; unnormalised
ANOVA-F statistics are dominated by high-amplitude delta-band features
regardless of their discriminative utility~\cite{Hastie2009}.

\subsection{Pipeline B: Extraction $\to$ Scaling $\to$ PCA}

\begin{algorithm}[h]
\caption{Pipeline B}
\label{alg:pipeB}
\begin{algorithmic}[1]
  \Require Raw feature matrix $\mathbf{X} \in \mathbb{R}^{m \times d}$
  \State $\mathbf{X}' \gets \texttt{MinMaxScaler}(\mathbf{X})$
    \Comment{bound to $[0,1]$, suppress spike artefacts}
  \State $\mathbf{X}'' \gets \texttt{StandardScaler}(\mathbf{X}')$
    \Comment{re-centre for PCA}
  \State $\mathbf{X}''' \gets \texttt{PCA}(\mathbf{X}'', n_\text{components}{=}10)$
    \Comment{project to principal subspace}
  \State \Return $\mathbf{X}'''$
\end{algorithmic}
\end{algorithm}

\noindent\textbf{Limitation:} PCA discards variance without regard for
class-discriminability---a significant disadvantage for minority-class
detection in severely imbalanced settings. PCA is known to discard
minority-class structure encoded in lower-variance principal
components~\cite{Hastie2009}.

% ============================================================
\section{Baseline Logistic Regression}
\label{sec:baseline}

The baseline model is L2-regularised logistic regression ($C{=}1.0$) trained
independently on the output of each pipeline, with a 20\% stratified
held-out test set. For a feature vector $\mathbf{x}$, the predicted seizure
probability is:
\begin{equation}
  P(y{=}1 \mid \mathbf{x}) = \frac{1}{1 + e^{-(\beta_0 + \boldsymbol{\beta}^\top\mathbf{x})}}.
  \label{eq:logreg}
\end{equation}

Three metrics are reported. \textbf{Accuracy} measures overall correctness
but inflates on imbalanced data. \textbf{F1-score} is the harmonic mean of
precision and recall. \textbf{PR-AUC} (area under the precision-recall
curve) is the \emph{primary metric} because it does not inflate from the
large true-negative count inherent in normal EEG epochs, unlike ROC-AUC
which remains high even when recall on the positive class is poor~\cite{Davis2006}.

\begin{table}[h]
\centering
\caption{Baseline Results ($C{=}1.0$, L2 Regularisation)}
\label{tab:baseline}
\begin{tabular}{llcccc}
\toprule
\textbf{Dataset} & \textbf{Pipe} &
\textbf{Acc} & \textbf{F1} & \textbf{PR-AUC} & \textbf{ROC-AUC} \\
\midrule
CHB-MIT  & A & 0.962 & 0.756 & 0.766 & 0.905 \\
CHB-MIT  & B & 0.927 & 0.425 & 0.554 & 0.854 \\
\midrule
Bonn-EEG & A & 0.960 & 0.902 & 0.937 & 0.954 \\
Bonn-EEG & B & 0.922 & 0.805 & 0.854 & 0.931 \\
\midrule
Kaggle   & A & 0.985 & 0.980 & 0.984 & 0.988 \\
Kaggle   & B & 0.980 & 0.974 & 0.977 & 0.983 \\
\bottomrule
\end{tabular}
\end{table}

\begin{figure}[htbp]
  \centering
  \includegraphics[width=\columnwidth]{fig2_pr_curves_baseline.png}
  \caption{Precision-Recall curves for baseline LR ($C{=}1.0$) under
           Pipeline~A (solid) and Pipeline~B (dashed). Pipeline~A achieves
           higher area on all three datasets, most notably on CHB-MIT
           ($\Delta{=}+0.212$).}
  \label{fig:pr_baseline}
\end{figure}

% ============================================================
\section{Overfitting and Underfitting Analysis}
\label{sec:overfitting}

Three intentional $C$-value scenarios are constructed on CHB-MIT to
demonstrate both failure modes of logistic regression.

\noindent\textbf{Underfitting ($C{=}0.0001$):} The penalty is so strong
that all coefficients collapse toward zero. The model predicts the majority
class exclusively ($\text{F1}{=}0.000$)---textbook excessive bias, with
training and validation scores coinciding at floor level.

\noindent\textbf{Balanced ($C{=}1.0$):} Training F1 ($0.804$) and test F1
($0.807$) are high with negligible gap ($-0.003$)---the bias-variance sweet
spot.

\noindent\textbf{Overfitting ($C{=}10{,}000$):} Regularisation is
effectively removed. The 5-fold cross-validation sweep
(Fig.~\ref{fig:overfit}, left) reveals a widening train/validation gap in
the high-$C$ zone, confirming excess variance.

\begin{table}[h]
\centering
\caption{Intentional Over/Underfitting Scenarios (CHB-MIT)}
\label{tab:scenarios}
\begin{tabular}{lcccc}
\toprule
\textbf{Scenario} & $C$ & \textbf{Train F1} & \textbf{Test F1} & \textbf{Gap} \\
\midrule
Underfitting & 0.0001    & 0.000 & 0.000 & $0.000$ \\
Balanced     & 1.0       & 0.804 & 0.807 & $-0.003$ \\
Overfitting  & $10{,}000$ & 0.802 & 0.815 & $-0.013$ \\
\bottomrule
\end{tabular}
\end{table}

\begin{figure}[htbp]
  \centering
  \includegraphics[width=\columnwidth]{fig3_overfit_underfit.png}
  \caption{Left: Validation curve (F1 vs.\ $C$, log-scale) with annotated
           underfitting and overfitting zones and shaded generalisation gap.
           Right: Learning curve showing train/validation convergence as
           training set size increases.}
  \label{fig:overfit}
\end{figure}

The learning curve (Fig.~\ref{fig:overfit}, right) shows that training and
validation F1 converge as data grows, with the gap narrowing to ${\approx}0.05$
at full size---consistent with a model of appropriate capacity that would
benefit from additional data.

% ============================================================
\section{Regularisation Study}
\label{sec:regularisation}

Three regularisation strategies are evaluated across an 11-point $C$ grid
$[0.001, \ldots, 100.0]$ using Pipeline~A features on all three datasets.

\subsection{L2 (Ridge)}
Penalises the squared coefficient norm. Shrinks all coefficients uniformly
but retains all features. The objective function is:
\begin{equation}
  J(\mathbf{W},b) = \frac{1}{m}\sum_{i=1}^{m}\mathcal{L}(\hat{y}^{(i)}, y^{(i)})
  + \frac{\lambda}{2m}\sum\|\mathbf{W}\|^2.
  \label{eq:l2}
\end{equation}

\subsection{L1 (Lasso)}
Penalises the absolute coefficient norm, driving some coefficients exactly to
zero---implicit feature selection. Particularly valuable for high-dimensional
EEG where many spectral features are uninformative.

\subsection{Elastic Net}
Combines L1 and L2 with mixing ratio $\rho{=}0.5$. Inherits L1's sparsity
and L2's grouping effect (correlated features receive similar non-zero
weights):
\begin{equation}
  J(\mathbf{W},b) = \frac{1}{m}\sum_{i=1}^{m}\mathcal{L}(\hat{y}^{(i)}, y^{(i)})
  + \alpha\rho\|\mathbf{W}\|_1
  + \frac{\alpha(1-\rho)}{2}\|\mathbf{W}\|^2.
  \label{eq:enet}
\end{equation}

\begin{table}[h]
\centering
\caption{Regularisation at $C{=}1.0$, Pipeline~A}
\label{tab:regularisation}
\begin{tabular}{llccc}
\toprule
\textbf{Dataset} & \textbf{Reg.} & \textbf{F1} & \textbf{PR-AUC} & \textbf{Sparsity} \\
\midrule
CHB-MIT  & L2           & 0.756 & 0.766 & 0.0\% \\
CHB-MIT  & L1           & 0.756 & 0.766 & 0.0\% \\
CHB-MIT  & Elastic Net  & 0.756 & 0.766 & 0.0\% \\
\midrule
Bonn-EEG & L2           & 0.902 & 0.937 & 0.0\% \\
Bonn-EEG & L1           & 0.902 & 0.938 & 6.7\% \\
Bonn-EEG & Elastic Net  & 0.902 & 0.937 & 0.0\% \\
\midrule
Kaggle   & L2           & 0.980 & 0.984 & 6.7\%  \\
Kaggle   & L1           & 0.980 & 0.984 & 26.7\% \\
Kaggle   & Elastic Net  & 0.980 & 0.984 & 20.0\% \\
\bottomrule
\end{tabular}
\end{table}

\begin{figure}[htbp]
  \centering
  \includegraphics[width=\columnwidth]{fig4_regularisation_study.png}
  \caption{F1 and PR-AUC vs.\ $C$ (log scale) for L1, L2, and Elastic Net
           across all three datasets. L1 degrades fastest at low $C$ due to
           coefficient collapse. All regularisers converge at $C{=}1.0$
           after upstream SelectKBest feature selection.}
  \label{fig:reg_curves}
\end{figure}

\begin{figure}[htbp]
  \centering
  \includegraphics[width=0.75\columnwidth]{fig5_sparsity_heatmap.png}
  \caption{Coefficient sparsity (\% $|w| < 10^{-4}$) at $C{=}1.0$.
           L1 achieves the highest sparsity (26.7\% on Kaggle);
           L2 produces no zero coefficients.}
  \label{fig:sparsity}
\end{figure}

\begin{table}[h]
\centering
\caption{Mean PR-AUC over Full $C$ Grid $[0.001 \ldots 100]$}
\label{tab:mean_prauc}
\begin{tabular}{lcccc}
\toprule
\textbf{Dataset} & \textbf{L2} & \textbf{Elastic Net} & \textbf{L1} &
$\Delta_\text{L2-L1}$ \\
\midrule
CHB-MIT  & \textbf{0.761} & 0.705 & 0.687 & $+0.074$ \\
Bonn-EEG & \textbf{0.934} & 0.867 & 0.840 & $+0.094$ \\
Kaggle   & \textbf{0.984} & 0.929 & 0.923 & $+0.061$ \\
\bottomrule
\end{tabular}
\end{table}

The $C$-sweep (Fig.~\ref{fig:reg_curves}) reveals the critical difference:
L1 degrades sharply at low $C$ because it nullifies all coefficients, leaving
only the intercept. Elastic Net maintains non-trivial performance at lower $C$
by retaining the L2 grouping term, while L2 is most stable throughout. After
Pipeline~A's SelectKBest step pre-applies sparsity, additional L1 sparsity is
counterproductive at low $C$.

\begin{figure}[htbp]
  \centering
  \includegraphics[width=0.9\columnwidth]{fig10_regulariser_heatmap.png}
  \caption{Mean PR-AUC heatmap over the full $C$ grid. L2 leads on all
           datasets; Elastic Net is second-best; L1 is weakest, especially
           on CHB-MIT ($\Delta{=}0.074$).}
  \label{fig:reg_heatmap}
\end{figure}

% ============================================================
\section{Class Imbalance Handling}
\label{sec:imbalance}

Class imbalance is the dominant practical challenge in seizure prediction.
Four strategies are compared at $C{=}1.0$ with Pipeline~A features:

\begin{enumerate}
  \item \textbf{No correction}---model biased toward majority class.
  \item \textbf{SMOTE}~\cite{Chawla2002}---generates synthetic seizure epochs
        by $k$-nearest neighbour interpolation in feature space.
  \item \textbf{Random undersampling}---discards majority-class samples until
        balanced; computationally cheap but discards real data.
  \item \textbf{Class weighting}---reweights the logistic loss by inverse
        class frequency, equivalent to oversampling without altering the
        training set.
\end{enumerate}

\begin{table}[h]
\centering
\caption{Imbalance Handling on CHB-MIT (9.2\% seizure, most severe)}
\label{tab:imbalance}
\begin{tabular}{lcccc}
\toprule
\textbf{Method} & \textbf{Prec.} & \textbf{Recall} & \textbf{F1} &
\textbf{PR-AUC} \\
\midrule
No correction   & \textbf{0.922} & 0.641 & \textbf{0.756} & 0.766 \\
SMOTE           & 0.468          & \textbf{0.783} & 0.585 & 0.761 \\
Undersampling   & 0.450          & \textbf{0.783} & 0.571 & 0.758 \\
Class weighting & 0.449          & 0.772          & 0.568 & \textbf{0.767} \\
\bottomrule
\end{tabular}
\end{table}

\begin{figure}[htbp]
  \centering
  \includegraphics[width=\columnwidth]{fig6a_imbalance_scatter.png}
  \caption{Precision vs.\ Recall scatter per imbalance method across all
           three datasets. Correction methods shift the operating point to
           higher recall at the cost of precision.}
  \label{fig:imbalance_scatter}
\end{figure}

\begin{figure}[htbp]
  \centering
  \includegraphics[width=\columnwidth]{fig6b_prauc_imbalance.png}
  \caption{PR-AUC by imbalance method across datasets. Differences are
           largest on CHB-MIT; on Kaggle (near-balanced) all methods
           converge.}
  \label{fig:imbalance_bar}
\end{figure}

A critical insight from Table~\ref{tab:imbalance}: no-correction achieves
the highest precision ($0.922$) and F1 ($0.756$) on CHB-MIT, but masks poor
recall ($0.641$)---only two-thirds of seizures are detected. In clinical
deployment, missed seizures (false negatives) are more dangerous than false
alarms. SMOTE and class weighting substantially improve recall ($0.783$ and
$0.772$) with comparable PR-AUC ($0.761$ and $0.767$), confirming they shift
the precision-recall operating point rather than improving the underlying
discriminator.

% ============================================================
\section{Comparative Analysis}
\label{sec:analysis}

This section addresses the four key research questions posed in the
assignment, synthesising evidence from Sections~\ref{sec:baseline}--\ref{sec:imbalance}.

\subsection{Q1: Does Preprocessing Order Affect Results?}

\textbf{Yes, significantly.} Pipeline~A outperforms Pipeline~B across all
three datasets (Table~\ref{tab:baseline}). The PR-AUC advantage is largest
on CHB-MIT ($+0.213$, from $0.554$ to $0.766$)---the dataset with the most
features ($64$) and worst class imbalance. Normalising before ANOVA-F
scoring is structurally necessary: unnormalised statistics are biased toward
high-variance features regardless of discriminative utility. PCA in
Pipeline~B compounds this by discarding minority-class variance encoded in
lower-variance principal components---a known failure mode of unsupervised
dimensionality reduction on imbalanced data~\cite{Hastie2009}.

\begin{figure}[htbp]
  \centering
  \includegraphics[width=\columnwidth]{fig9_pipeline_summary.png}
  \caption{F1 and PR-AUC by dataset and pipeline for baseline L2.
           Pipeline~A dominates on both metrics across all three datasets.}
  \label{fig:pipeline_summary}
\end{figure}

\subsection{Q2: Which Regulariser Generalises Best?}

\textbf{L2 (Ridge).} It achieves the highest mean PR-AUC over the $C$ grid
on all three datasets (Table~\ref{tab:mean_prauc}: $0.761$, $0.934$,
$0.984$). Elastic Net is second-best; L1 is consistently weakest, with the
largest gap on CHB-MIT (L2$-$L1$=0.074$ PR-AUC). L2's stability arises
from its quadratic penalty, which degrades smoothly across the $C$ range
rather than producing abrupt coefficient zeroing. With $k{=}15$ features
already selected by SelectKBest, additional L1 sparsity is counterproductive
at low $C$.

\subsection{Q3: Does Elastic Net Consistently Outperform L1 and L2?}

\textbf{Elastic Net outperforms L1 in all cases, but does not consistently
outperform L2} in this study. This nuanced finding implies that Elastic
Net's advantage is most pronounced in high-dimensional settings
($d > 30$ post-preprocessing features) where L1's coefficient collapse at
low $C$ is a real risk. Here, Pipeline~A's SelectKBest step pre-applies
sparsity, reducing the need for L1's within-model feature selection and
making L2 sufficient. The practical implication: Elastic Net is the safer
default when no upstream feature selection is performed.

\begin{figure}[htbp]
  \centering
  \includegraphics[width=0.85\columnwidth]{fig8_radar_regulariser.png}
  \caption{Radar chart comparing regularisers across five metrics on
           Bonn-EEG ($C{=}1.0$, Pipeline~A). L2 and Elastic Net achieve
           near-identical multi-metric profiles; L1 lags on coefficient
           density.}
  \label{fig:radar}
\end{figure}

\subsection{Q4: How Does Imbalance Handling Interact with Regularisation?}

The interaction effect is strongest on CHB-MIT ($9.2\%$).
Fig.~\ref{fig:interaction} presents the full PR-AUC interaction matrix.
Class weighting $+$ L2 ($0.767$) and SMOTE $+$ L2 ($0.761$) produce the
highest PR-AUC; differences are within experimental variance. Any correction
method is preferable to no correction when recall is the primary concern.
The interaction is essentially absent on Kaggle (near-balanced), confirming
that imbalance handling is only necessary when the class ratio exceeds
approximately $3:1$.

\begin{figure}[htbp]
  \centering
  \includegraphics[width=\columnwidth]{fig7_interaction_heatmap.png}
  \caption{PR-AUC interaction heatmap: imbalance method $\times$ regulariser
           on CHB-MIT ($C{=}1.0$, Pipeline~A). Class weighting combinations
           yield the highest PR-AUC ($0.767$).}
  \label{fig:interaction}
\end{figure}

% ============================================================
\section{Conclusion}
\label{sec:conclusion}

This study makes three principal contributions.

First, it provides quantitative evidence that \emph{preprocessing order is
not a stylistic choice}: StandardScaler before SelectKBest yields $+0.213$
PR-AUC on CHB-MIT compared to the PCA-based pipeline. Normalising before
feature ranking is a structural necessity for scale-invariant ANOVA-F
scoring.

Second, it characterises the $C$-sweep behaviour of L1, L2, and Elastic Net,
showing that \emph{L2 achieves the best mean PR-AUC when combined with
upstream feature selection}, and that \emph{Elastic Net is the safer default
in its absence}. Elastic Net consistently outperforms L1 by avoiding
coefficient collapse at low $C$.

Third, it demonstrates that \emph{class-imbalance correction methods shift
the precision-recall operating point rather than improving PR-AUC
substantially}---motivating PR-AUC as the primary selection criterion over
F1 or accuracy for clinical seizure detection tasks.

Future work should extend this framework to non-linear models (SVM with RBF
kernel, gradient boosting) and to cross-patient generalisation scenarios
where the covariate shift between training and test distributions is more
severe than the within-patient setting studied here.

% ============================================================
% Balance last two columns on last page
\balance

\begin{thebibliography}{9}

\bibitem{WHO2019}
World Health Organization,
\textit{Epilepsy: A Public Health Imperative},
WHO Report, 2019.

\bibitem{Goldberger2010}
A.~L.~Goldberger \textit{et al.},
``CHB-MIT Scalp EEG Database,''
\textit{PhysioNet}, 2010.
doi:10.13026/C2K01R.

\bibitem{Andrzejak2001}
R.~G.~Andrzejak \textit{et al.},
``Indications of nonlinear deterministic and finite-dimensional structures in
time series of brain electrical activity,''
\textit{Physical Review E}, vol.~64, no.~6, 2001.

\bibitem{Kaggle}
Kaggle Seizure Prediction Challenge Dataset.
[Online]. Available: \url{https://www.kaggle.com/c/seizure-prediction}

\bibitem{Hastie2009}
T.~Hastie, R.~Tibshirani, and J.~Friedman,
\textit{The Elements of Statistical Learning},
2nd~ed., Springer, 2009, ch.~3.

\bibitem{Davis2006}
J.~Davis and M.~Goadrich,
``The relationship between precision-recall and ROC curves,''
in \textit{Proc.\ ICML}, pp.~233--240, 2006.

\bibitem{Chawla2002}
N.~V.~Chawla, K.~W.~Bowyer, L.~O.~Hall, and W.~P.~Kegelmeyer,
``SMOTE: Synthetic minority over-sampling technique,''
\textit{JAIR}, vol.~16, pp.~321--357, 2002.

\bibitem{Pedregosa2011}
F.~Pedregosa \textit{et al.},
``Scikit-learn: Machine learning in Python,''
\textit{JMLR}, vol.~12, pp.~2825--2830, 2011.

\bibitem{Lemaitre2017}
G.~Lema\^{i}tre, F.~Nogueira, and C.~K.~Aridas,
``Imbalanced-learn: A python toolbox to tackle the curse of imbalanced
datasets in machine learning,''
\textit{JMLR}, vol.~18, no.~17, pp.~1--5, 2017.

\end{thebibliography}

\end{document}
